\section{Goal-generated polynomial relation closure}
\label{sec:vector}

\subsection{A small problem class with a useful infinite quantifier}

Let $A$ be a finite nonempty alphabet. For each $a\in A$, let $F_a:\Q^d\to\Q^d$ be a total polynomial map with rational coefficients. Fix $s_0\in\Q^d$ and target polynomials $g_1,\ldots,g_t\in\Q[X_1,\ldots,X_d]$. A word $w=a_1\cdots a_\ell$ acts in execution order:
\[
 F_w=F_{a_\ell}\circ\cdots\circ F_{a_1},\qquad F_\varepsilon=\mathrm{id}.
\]
The requested assertion is
\begin{equation}
 \forall w\in A^*,\quad \forall j\le t,\qquad g_j(F_w(s_0))=0.
 \label{eq:allwords}
\end{equation}
One symbol describes a recurrence iterated an arbitrary number of times. Several symbols describe every interleaving of a fixed collection of updates. Running two implementations on the same word is represented by a product state, so the targets may express relations between them.

This input class does not include guards, partial transitions, truncated subtraction, integer division, or arbitrary recursive data. A frontend may reduce a larger problem to this class only with a proved bridge. In particular, a negative execution of an unguarded abstraction need not be a feasible execution of a guarded source program.

\begin{definition}[Pullback]
For a polynomial $q$, define $F_a^*q=q\circ F_a$. This is substitution of the coordinate polynomials of $F_a$ for the variables of $q$, performed \emph{simultaneously}.
\end{definition}
A pullback asks which property before a step would imply the desired property after the step. Unlike forward trajectory sampling, it preserves a symbolic statement about every state.

\subsection{The certificate and its complete soundness argument}

Search for a column $B=(b_1,\ldots,b_r)^T$ of rational polynomials, matrices $H_a\in\Q^{r\times r}$, and a matrix $C\in\Q^{t\times r}$ satisfying
\begin{align}
 B(s_0)&=0, \label{eq:vector-init}\\
 B(F_a(s))&=H_a B(s) &&(a\in A),\label{eq:vector-step}\\
 (g_1(s),\ldots,g_t(s))^T&=C B(s).\label{eq:vector-target}
\end{align}
The second and third lines are polynomial identities, not equalities tested on a finite grid. Individual $b_i$ may grow, alternate sign, or mix with other entries. Only their common zero set must be preserved.

\begin{theorem}[Vector closure certificate]
If \eqref{eq:vector-init}--\eqref{eq:vector-target} hold, then \eqref{eq:allwords} holds.
\end{theorem}
\begin{proof}
Induct on execution length. The empty execution has $B(s_0)=0$. If $B(s)=0$ and the next symbol is $a$, then $B(F_a(s))=H_a0=0$. Hence every reachable state has $B=0$. Multiplication by $C$ gives every target equality.
\end{proof}

The checker does not need to verify that the $b_i$ are independent, minimal, or actually discovered by the advertised algorithm. Those properties affect search cost, not soundness. It checks the original targets, every original transition, the initial evaluations, and the displayed identities. The search-origin annotations are deliberately irrelevant to positive acceptance.

\subsection{Discovering the basis from the target}

Begin with the targets in a FIFO worklist and an empty basis. For each polynomial $q$ removed from the worklist, test whether it belongs to the rational span of the current basis. If so, discard it. Otherwise append it and enqueue its pullbacks along every transition. Exact row reduction implements span membership.

\begin{lstlisting}
B := []; work := [(g_j, target=j, word=[]) for j]
while work is nonempty:
    (q, j, w) := pop_front(work)
    if polynomial or work budgets exceeded: return UNKNOWN
    if q(s0) != 0: return checked_word_counterexample(j, w)
    if q is in span_Q(B): continue
    if basis budget exhausted: return UNKNOWN
    append q to B
    for a in alphabet:
        enqueue(q composed with F_a, j, [a] ++ w)
reconstruct H_a and C over exact rationals
return certificate only after independent replay
\end{lstlisting}

This is not enumeration of all degree-$D$ invariant templates. Only the closure generated by the goal is explored. It resembles a block Krylov construction, but the operators are symbolic pullbacks and there may be several of them. The implementation recomputes exact coordinate systems for clarity; an incremental echelon basis is a straightforward performance improvement.

\paragraph{Why a discarded polynomial needs no descendants.}
If $q=\sum_i c_i b_i$, then $q\circ F_a=\sum_i c_i(b_i\circ F_a)$. Once the basis descendants are accounted for, the descendants of $q$ add no new direction. This is the mathematical justification for the worklist pruning; it does not depend on a heuristic estimate of relevance.

\paragraph{Why word order is easy to get wrong.}
Suppose $q=g_j\circ F_w$. Then
\[
 q\circ F_a=g_j\circ F_w\circ F_a=g_j\circ F_{aw}.
\]
Pulling back along $a$ \emph{prepends} $a$ to the execution word. The implementation stores this exact provenance. A test with updates $F_0(x,y)=(1,y)$ and $F_1(x,y)=(x,x)$, initial state $(0,0)$, and target $y=0$ finds $[0,1]$. Its reverse $[1,0]$ is not a counterexample, and the independent checker rejects that mutation.

\subsection{Termination and relative completeness}

Let
\[
 V=\Span_\Q\{g_j\circ F_w:j\le t,\ w\in A^*\}.
\]

\begin{theorem}[Finite orbit-span guarantee]
If $\dim V=r<\infty$, exact uncapped vector closure terminates. If \eqref{eq:allwords} is true, it returns a positive certificate. If \eqref{eq:allwords} is false, it returns a concrete counterexample. At most $r$ basis insertions and $|A|r$ newly generated pullbacks are needed.
\end{theorem}
\begin{proof}
Every inserted polynomial lies in $V$ and is independent of the previous ones, so there are at most $r$ insertions. Each insertion enqueues finitely many descendants. The queue is therefore finite. At exhaustion, the resulting span contains all targets and is invariant under every pullback, because each basis pullback was processed. It follows by induction on word length that it equals $V$.

If some processed $q=g_j\circ F_w$ has $q(s_0)\ne0$, its stored word refutes the assertion. If no such polynomial exists, every basis member vanishes at $s_0$, giving the initial certificate condition. Exact coordinate reconstruction gives the other conditions. Were there an unreported counterexample after closure, its target pullback would be a linear combination of basis polynomials, all zero at $s_0$, a contradiction.
\end{proof}

\begin{corollary}[Affine transition fragment]
For affine $F_a$ and targets of total degree at most $D$, vector closure terminates in a space of dimension at most $\binom{d+D}{D}$.
\end{corollary}
\begin{proof}
Affine substitution preserves the degree bound. The monomials of total degree at most $D$ form a vector-space basis of the stated size.
\end{proof}

The dimension bound does not make exact arithmetic cheap. Rational numerators and denominators can grow, and a naive implementation repeatedly builds dense matrices over all current monomials. The prototype has explicit degree, term, basis, and pullback ceilings. A ceiling gives \unknown{}, even on a mathematically complete fragment. The counterexample is not advertised as shortest in this algebraic worker: linear pruning and multiple targets complicate that additional promise. Only the finite-machine BFS below promises shortest distinguishing words.

\subsection{Worked example: a coupled quadratic relation}

Consider synchronized Fibonacci-style updates
\[
 F(x,y,u,v)=(y,x+y,v,u+v),\qquad s_0=(0,1,0,1),
\]
and the target
\[
 b_0=x^2+y^2-u^2-v^2.
\]
The target is not conserved as a polynomial. Its pullback is
\[
 b_1=x^2+2xy+2y^2-u^2-2uv-2v^2.
\]
A second pullback satisfies
\[
 b_1\circ F=-b_0+3b_1.
\]
Thus the search discovers
\begin{equation}
 B=\begin{pmatrix}b_0\\b_1\end{pmatrix},\qquad
 H=\begin{pmatrix}0&1\\-1&3\end{pmatrix},\qquad C=\begin{pmatrix}1&0\end{pmatrix}.
 \label{eq:coupled-example}
\end{equation}
Both entries vanish initially, so the joint relation is inductive. There is no need to guess $b_1$ in advance, enumerate quadratic invariants, or ask a nonlinear solver to invent a matrix and basis simultaneously.

The particular synchronized example has an easier human proof by equality of the two state pairs. Its purpose is not to make a benchmark look difficult. It isolates the discovery mechanism and exposes precisely why requiring $b_0\circ F=b_0$ would be too restrictive. The generated affine corpus varies the two updates and the quadratic observation rather than claiming this single example measures realistic Lean difficulty.

\subsection{Worked example: a sign-changing invariant}

With
\[
 F(x,y,t)=(y,x+y,-t),\quad s_0=(0,1,1),\quad
 b=y^2-xy-x^2-t,
\]
exact expansion gives $b\circ F=-b$. The worker returns the one-by-one matrix $H=(-1)$. The conserved-quantity condition $b\circ F=b$ fails, but the zero set remains invariant. This gives a Cassini-shaped relation while representing the alternating sign by an extra state coordinate, rather than incorrectly reifying $(-1)^n$ as a polynomial in $n$.

\subsection{A precise obstruction to vector closure}

For $F(x,y)=(x^2,y^2)$ and $g=x-y$, the successive pullbacks are $x^{2^j}-y^{2^j}$. They are linearly independent over $\Q$: in a finite combination, the greatest exponent occurs in exactly one summand and cannot cancel. Hence $V$ is infinite-dimensional. An arbitrarily generous finite vector-space budget cannot establish closure this way.

Nevertheless, the invariant is elementary:
\[
 (x-y)\circ F=(x+y)(x-y).
\]
The obstacle is the restriction to \emph{constant} coefficients. Allowing a polynomial multiplier changes the mathematics of closure and motivates the next worker. This is an algorithmic routing criterion, not merely a request to increase fuel.
